% GED Math: Number Sense and Problem Solving Lessons
% LaTeX Beamer Slide Deck

\documentclass{beamer}
\usetheme{Madrid}
\usepackage{amsmath, amssymb}
\title{GED Math: Number Sense and Problem Solving}

\date{\today}

\begin{document}

\frame{\titlepage}

%==================================================
% Lesson 1: Understanding Numbers
%==================================================

\begin{frame}{Lesson 1: Understanding Numbers}
\begin{itemize}
    \item Types of numbers: natural, whole, integers, rational, irrational.
    \item Prime and composite numbers.
    \item Factors and multiples.
    \item Greatest Common Factor (GCF) and Least Common Multiple (LCM).
\end{itemize}
\end{frame}

\begin{frame}{Example 1: Classifying Numbers}
\textbf{Example:} Classify the number $\sqrt{16}$.

\textbf{Solution:}
\begin{itemize}
    \item $\sqrt{16} = 4$, which is a whole number.
    \item Therefore, it is natural, whole, integer, and rational.
\end{itemize}
\end{frame}

\begin{frame}{Example 2: Prime Factorization}
\textbf{Example:} Find the prime factorization of 84.

\textbf{Solution:}
\begin{align*}
84 &= 2 \times 42 = 2 \times 2 \times 21 = 2^2 \times 3 \times 7
\end{align*}
\end{frame}

%==================================================
% Lesson 2: Fractions and Decimals
%==================================================

\begin{frame}{Lesson 2: Fractions and Decimals}
\begin{itemize}
    \item Converting between fractions, decimals, and percentages.
    \item Simplifying fractions.
    \item Adding, subtracting, multiplying, and dividing fractions.
\end{itemize}
\end{frame}

\begin{frame}{Example 3: Simplifying a Fraction}
Simplify $\frac{18}{24}$.

\textbf{Solution:}
\begin{align*}
\text{GCF}(18,24) = 6, \quad \frac{18}{24} = \frac{3}{4}
\end{align*}
\end{frame}

\begin{frame}{Example 4: Fraction to Decimal}
Convert $\frac{7}{8}$ to a decimal.

\textbf{Solution:}
\begin{align*}
\frac{7}{8} = 0.875
\end{align*}
\end{frame}

%==================================================
% Lesson 3: Ratios, Rates, and Proportions
%==================================================

\begin{frame}{Lesson 3: Ratios, Rates, and Proportions}
\begin{itemize}
    \item Understanding ratios and rates.
    \item Solving proportions.
    \item Applications in real-world problems.
\end{itemize}
\end{frame}

\begin{frame}{Example 5: Solving a Proportion}
Solve $\frac{3}{4} = \frac{x}{12}$.

\textbf{Solution:}
\begin{align*}
3 \times 12 = 4x \Rightarrow x = 9
\end{align*}
\end{frame}

%==================================================
% Lesson 4: Percents
%==================================================

\begin{frame}{Lesson 4: Percents}
\begin{itemize}
    \item Converting between fractions, decimals, and percents.
    \item Finding percentage of a number.
    \item Percentage increase and decrease.
\end{itemize}
\end{frame}

\begin{frame}{Example 6: Finding a Percent}
Find 25\% of 80.

\textbf{Solution:} $0.25 \times 80 = 20$
\end{frame}

\begin{frame}{Example 7: Percentage Change}
The price increases from 50 to 65. Find the percent increase.

\textbf{Solution:}
\begin{align*}
\text{Percent change} &= \frac{65 - 50}{50} \times 100\% = 30\%
\end{align*}
\end{frame}

%==================================================
% Lesson 5: Estimation and Rounding
%==================================================

\begin{frame}{Lesson 5: Estimation and Rounding}
\begin{itemize}
    \item Rounding to the nearest whole number, tenth, hundredth.
    \item Estimation in problem solving.
\end{itemize}
\end{frame}

\begin{frame}{Example 8: Rounding}
Round 67.489 to the nearest tenth.

\textbf{Solution:} 67.489 $\approx$ 67.5
\end{frame}

%==================================================
% Lesson 6: Order of Operations
%==================================================

\begin{frame}{Lesson 6: Order of Operations}
\begin{itemize}
    \item Use PEMDAS: Parentheses, Exponents, Multiplication/Division, Addition/Subtraction.
\end{itemize}
\end{frame}

\begin{frame}{Example 9: Evaluating an Expression}
Evaluate $6 + 3 \times (2 + 4)^2$.

\textbf{Solution:}
\begin{align*}
6 + 3 \times 6^2 &= 6 + 3 \times 36 = 6 + 108 = 114
\end{align*}
\end{frame}

%==================================================
% Lesson 7: Absolute Value and Number Line
%==================================================

\begin{frame}{Lesson 7: Absolute Value}
\begin{itemize}
    \item $|x|$ is the distance from zero.
    \item Example: $|-5| = 5$
\end{itemize}
\end{frame}

\begin{frame}{Example 10: Comparing Numbers}
Which is greater: $|-3|$ or $2$?

\textbf{Solution:} $|-3| = 3$, so $3 > 2$.
\end{frame}

%==================================================
% Continue examples up to 25
%==================================================

\begin{frame}{Example 25: Multi-step Problem}
A car travels 150 miles in 3 hours. How far will it travel in 5 hours at the same rate?

\textbf{Solution:}
\begin{align*}
\text{Rate} &= 150/3 = 50 \text{ mph} \\
\text{Distance} &= 50 \times 5 = 250 \text{ miles}
\end{align*}
\end{frame}

\end{document}
